\documentclass[output=paper]{langscibook}
\ChapterDOI{10.5281/zenodo.22880258}
\title{\textit{A sinoo na inu} – The frame of the house}
\author{Enoch Horai Magum}
\abstract{}
\begin{document}
\maketitle\label{chap:16}
\begin{figure}[H]
    \includegraphics[trim=0 0 0 0,clip,width=\linewidth]{figures/16_1.png}
    \caption{\textit{A sinoo na inu} -- The timber frame (lit. `skeleton') of the house}
    \label{fig:16_1}
\end{figure}
\begin{xltabular}{\textwidth}{lXX}
    (\ref{ex:text11-1})&\textit{Amaa humhum na inu gorogoroho vaa Teapu}&Some parts of the Teop sleeping house\\
    (\ref{ex:text11-2})&\textit{\textbf{Toon}}&\textbf{Post}\\
    (\ref{ex:text11-3})&\textit{A ta naono vai a ta beera, bebeahu ae kikis to toon riori bona inu,}&A big, long and strong pole that is used as a post of the house,\\
    (\ref{ex:text11-4})&\textit{ge eve o naono to dee maana bona inu.}&or, (put differently,) it is a pole that carries the house.\\
    (\ref{ex:text11-5})&\textit{Bea ta tovihi, na tei koana tea beera na inu.}&How many, that depends on the size of the house.\\
    (\ref{ex:text11-6})&\textit{Evehee sa gonogono vahogo haari,}&But they do not get them in a careless way,\\
    (\ref{ex:text11-7})&\textit{ahiki, o naono to mataa ki nana bono toon o toraha.}&no, the tree that is good for a post is the New Guinea Teak.\\
    (\ref{ex:text11-8})&\textit{O naono o kikis ae o vaanagu koava.}&The wood is strong and resistant to rotting.\\
\end{xltabular}
\begin{xltabular}{\textwidth}{lXX}
    (\ref{ex:text11-9})&\textit{\textbf{Paee}}&\textbf{Horizontal beam} (joist, tie beam, and the main ridge pole)\\
    (\ref{ex:text11-10})&\textit{O naono o bebeahu to vaobete riori toon na maa paos.}&The long beam that is put on top of the posts.\\
    (\ref{ex:text11-11})&\textit{Evehee toro naono vai haabana o kikis, tabae eve to kahi dee nana bono paranda.}&But it must again be a strong beam, because it will carry the floor.\\
    (\ref{ex:text11-12})&\textit{A ta tovihi? Na tei koana tea bero tea toon.}&How many? That depends on the number of the posts.\\
    (\ref{ex:text11-13})&\textit{O hum to teitei bata nae mee, ei toon na maa sunsun.}&A place where they also lie, this is on top of the posts called \textit{sunsun} (the standing posts).\\
\end{xltabular}
\begin{xltabular}{\textwidth}{lXX}
    (\ref{ex:text11-14})&\textit{\textbf{Sunsun, suusun}}&\textbf{Standing post}\\
    (\ref{ex:text11-15})&\textit{O sunsun o naono to sun nana toon na maa paee.}&The standing post is a pole that stands on the horizontal beams.\\
    (\ref{ex:text11-16})&\textit{O meho paee repaa obete taono vaha ni maa bona,}&And another horizontal beam lies on it.\\
    (\ref{ex:text11-17})&\textit{O varuvaru repaa obete maa toon nae.}&And then the rafter lies on top of it (the horizonal beam).\\
    (\ref{ex:text11-18})&\textit{Eve o naono to kahi ato ni nana bono hum inu vaa gina.}&It is the pole that holds the upper part of the house.\\
\end{xltabular}
\begin{xltabular}{\textwidth}{lXX}
    (\ref{ex:text11-19})&\textit{\textbf{Gaasu}}&\textbf{King post}\\
    (\ref{ex:text11-20})&\textit{O gaasuu o naono to vavaasun riori tobina inu toon no paee.}&The kingpost is the pole that they stand in the middle of the house on top of  the tie beam.\\
    (\ref{ex:text11-21})&\textit{Beo tovihi? Na tei koana tea bebeahu na inu.}&How many? That just depends on the length of the house.\\
    (\ref{ex:text11-22})&\textit{Evehee a kiu tea maa naono bari mene dee}&The function of these poles is to carry (the main ridge pole)\\
    (\ref{ex:text11-23})&\textit{ge eve to vaatoka maa riori bono peho paee.}&or, (in other words,) it is the one on which they make the main ridge pole perch.\\
    (\ref{ex:text11-24})&\textit{Amaa ta naono bari mee amaa ta kakato koa.}&These poles (lit. `pieces of wood') are also quite short.\\
\end{xltabular}
\begin{xltabular}{\textwidth}{lXX}
    (\ref{ex:text11-25})&\textit{\textbf{Varuvaru}}&\textbf{Rafter}\\
    (\ref{ex:text11-26})&\textit{A varuvaru eve a naono to kahi vaobete mau ori toon no paee vai}&The rafters are the beams that are laid on the main ridgepole (lit. `horizontal beam')\\
    (\ref{ex:text11-27})&\textit{to paa vaobete vori toon na mata gaasuu.}&that is laid on the kingposts.\\
    (\ref{ex:text11-28})&\textit{Eve mee a naono to kahi vaave taono riori bona atovo.}&They are also the beams on which the sago palm thatches will be fixed.\\
    (\ref{ex:text11-29})&\textit{A bero nae mee na tei habana koana tea bebeahu na inu.}&Their number depends again only on the length of the house.\\
    (\ref{ex:text11-30})&\textit{A naono nae a vahavahara naono vaa tea apao.}&The trees used for them are little trees from the old garden.\\
    (\ref{ex:text11-31})&\textit{A varuvaru he toro tapaku tea buaku a papana tea papa matau ae papa keruka,}&The rafters must be made on the two sides, on the right side and the left side,\\
    (\ref{ex:text11-32})&\textit{ore paa vatasu vatatana tobina toon no paee vai}&so that they touch each other in the middle (of the roof) on top of the main ridgepole\\
    (\ref{ex:text11-33})&\textit{to paa vaatoka vori toon na maa gaasuu vaa tobina.}&that they have placed on the kingposts in the middle.\\
\end{xltabular}
\begin{xltabular}{\textwidth}{lXX}
    (\ref{ex:text11-34})&\textit{\textbf{Vaatetekoio}}&\textbf{Upper ridgepole}\\
    (\ref{ex:text11-35})&\textit{O naono bebeahu to kahi vaatoka vaha maa riori toon na maa varuvaru vai to paa vatasu vatatana vuru.}&A long beam that will be placed on top of the rafters that touch each other (on the lower ridgepole).\\
    (\ref{ex:text11-36})&\textit{Eve to kahi ato vakikis habana nana tea inu.}&It will also hold the house firmly together.\\
    (\ref{ex:text11-37})&\textit{Evehee a kiu teo naono bona rakaha, ei meno kovuaka.}&But the main function of this beam, this is for the top thatch sheet.\\
    (\ref{ex:text11-38})&\textit{Eve to kahi vaobete maa riori bono kovuaka ore paa obete vamataa koa maa teebona.}&It is the beam where they put the top thatch sheet so that it properly lies there.\\
\end{xltabular}
\begin{xltabular}{\textwidth}{lXX}
    (\ref{ex:text11-39})&\textit{\textbf{Hogo}}&\textbf{Floor joist}\\
    (\ref{ex:text11-40})&\textit{O naono to hogo vaha ni maa riori paana, be obete bata maana toon na maa paee vai tea maa toon.}&The beams that they put underneath (the floor), while they lie on the horizontal beams (that lie) on the posts.\\
    (\ref{ex:text11-41})&\textit{Eve to kahi obete mau o paranda.}&They are (the ones) on which the floor will lie.\\
    (\ref{ex:text11-42})&\textit{O tovihi? Na tei pete koana tea bebeahu ae tea tatabonata na inu.}&How many? That also just depends on the length and the width of the house.\\
\end{xltabular}
\begin{figure}[H]
    \includegraphics[width=\linewidth]{figures/16_2.jpg}
    \caption{\textit{A sinoo na inu} -- The frame of a house}
    \label{fig:16_2}
\end{figure}
\end{document}